\section{Síntesis y ampliación}

\subsection{Repaso del marco general de la clase}
Esta clase enlazó tres líneas: la línea de evaluación (qué métricas vale la pena optimizar), la línea de sistemas (cómo diseñar el harness y las herramientas) y la línea de seguridad (por qué es necesario verificar de forma dinámica). Las tres líneas convergen finalmente en la misma conclusión: \textbf{la competitividad de un Agent proviene de la calidad del ciclo cerrado del sistema, no de una capacidad textual de un solo intento}.

\begin{table}[H]
\centering
\begin{tabular}{p{3.1cm}p{4.4cm}p{4.2cm}p{2.6cm}}
\toprule
\textbf{Tema} & \textbf{Conclusión de esta clase} & \textbf{Práctica aplicable} & \textbf{Desafío pendiente} \\
\midrule
Diseño de evaluación & La evaluación orienta el rumbo de la investigación y el desarrollo, pero tiene una vida útil & Actualizar continuamente la distribución de tareas, calibrar con escenarios reales & Resistir la fuga de datos y el sobreajuste \\
Arquitectura del Agent & ReAct más retroalimentación de herramientas es la unidad básica & Diseño del flujo de control por capas, las trayectorias fallidas deben auditarse & Equilibrio entre costo y estabilidad \\
Ingeniería del harness & La interfaz de herramientas debe diseñarse junto con la capacidad del modelo & Pensamiento ACI, acciones compactas, salidas comprimibles & Compromiso entre generalidad y especialización \\
Tareas de seguridad & Se debe pasar del juicio estático a la verificación dinámica & sanitizer/debugger entran al bucle principal & Generalización entre proyectos y control de falsos positivos \\
Práctica de Big Sleep & El descubrimiento de vulnerabilidades reales ya es alcanzable & Iterar usando la evidencia de activación como función objetivo & Extenderse a pilas de sistemas más complejas \\
\bottomrule
\end{tabular}
\caption{Tabla resumen de las conclusiones centrales de Lecture 08}
\end{table}

\subsection{Further Reading}
\begin{itemize}
    \item Las líneas base centrales mencionadas en la conferencia de Sutton: MBPP, HumanEval, SWE-Bench, SWE-agent, AutoCodeRover.
    \item Evaluación y datos de seguridad: bases de datos públicas de vulnerabilidades (el ecosistema CVE), benchmarks de CTF, la línea de evaluación de seguridad de Meta, competencias relacionadas con DARPA.
    \item Dirección de ingeniería: Agent Computer Interface, análisis de fallos a nivel de trayectoria, diseño de harness que prioriza la verificación dinámica.
    \item Práctica de investigación en seguridad: materiales de divulgación pública de Big Sleep / Project Naptime, así como el proceso de divulgación de vulnerabilidades en colaboración con Project Zero.
\end{itemize}

\subsection{Resumen de la sección}
El juicio de Sutton en la etapa final puede tomarse directamente como hoja de ruta: \textit{``This is a wide open area. We're just getting started.''} Para los investigadores y los equipos de ingeniería, el siguiente punto de énfasis no es hacer otra ronda de puntuación superficial en benchmarks, sino integrar al Agent en flujos de trabajo de seguridad más reales, más auditables y más reproducibles.

\end{document}
